%%%%%%%%%%%%%%%%%%%%%%% CHAPTER - 3 %%%%%%%%%%%%%%%%%%%%%%%%%%
\chapter{A Review on Long-baseline Neutrino Oscillation Experiments}
\label{C3} 
%%%%%%%%%%%%%%%%%%%%%%%%%%%%%%%%%%%%%%%%%%%%%%%%%%%%%%%%%%%%%%
%%%%%%%%%%%%%%%%%%%%%%%%%%%%
Neutrinos, originating from various sources, differ not only in type but also in energy and the distance that they travel. As mentioned earlier, they are observed across an energy spectrum ranging from eV to PeV, depending on their sources and production mechanisms. For certain sources, both the energy of neutrinos and their propagation distance span a wide range (e.g., atmospheric and astrophysical neutrinos). Due to matter effects, neutrino mass-mixing parameters get evolved and the oscillation probabilities get modified in a non-trivial fashion as a function of energy and path length. Since these oscillation parameters $L$ and $E$ are inherently correlated, achieving high precision for a single parameter is challenging without addressing the uncertainties in the others. Therefore, a well-controlled, high-precision neutrino beam can be useful to tackle the uncertainties in neutrino energy and flux. This motivated the development of artificial setups designed to generate approximately mono-energetic neutrino beams, with a well-determined path from the source to the detector \cite{Agarwalla:2014fva}. \\
\par In 1960, Melvin Schwartz \cite{Schwartz:1960hg} proposed that an intense neutrino beam could be created using an accelerator, which would allow controlled laboratory-based studies of neutrino interactions. This idea was realized in the famous Brookhaven Neutrino Experiment (BNL E-734, 1962-1963) at the Alternating Gradient Synchrotron (AGS) \cite{Shiltsev:2019rfl}. The concept of producing a neutrino beam using accelerated protons became the foundation for modern long-baseline neutrino experiments. The experiment at BNL successfully demonstrated the existence of muon neutrino ($\nu_\mu$) as distinct from the electron neutrino ($\nu_e$). High-energy protons were directed onto a beryllium target, producing a shower of secondary particles, mainly pions and some kaons ($\pi^\pm$ and $K^\pm$). The pions and kaons decayed in a 200-meter-long decay tunnel, producing a beam of neutrinos via the decay process,
\begin{align}
    \pi^+ \rightarrow \mu^++\nu_\mu,\nonumber\\
    K^+ \rightarrow \mu^++\nu_\mu.
\end{align}
Since neutrinos feebly interact with matter, they pass through a thick steel shield ($\sim$13.5 m of steel), which absorbs all remaining hadrons and muons, leaving only a pure beam of muon neutrinos. A 10-ton aluminum spark chamber detector was placed 21 meters from the target to observe neutrino interactions. The neutrinos were expected to interact with nuclei in the aluminum via charged-current interaction,
\begin{align}
    \nu_\mu+n \rightarrow \mu^-+p.
\end{align}
The experiment only observed muons, proving that the beam contained a new neutrino species — the muon neutrino ($\nu_\mu$), distinct from the electron neutrino ($\nu_e$). This discovery was crucial in establishing the two-neutrino model, and in 1988, the Nobel Prize was awarded to Leon Lederman, Melvin Schwartz, and Jack Steinberger for this work \cite{Nobel1988}. In this chapter, we explore the key features of long-baseline experiments, outlining their general setups, primary objectives in advancing neutrino oscillation research, and significant contributions from past, present, and future \cite{Agarwalla:2024aee} LBL experiments worldwide.\\
%%%%%%%%%%%%%%%%%%%%%%%%%%%%%%%%%%%%%%%%%%%%%%%%%%%%%%%%%%%%%%%%%%%%%%%%%%%%%%%%%%%%%%%%%%%%%%%%%%%%%%%%%%%%%%%%%%%%%%%%%%%%%%%%%%%%%%%%%%%%%%%%%%%%%%%%%%%%%%%%%%%%%%%%%%%%
\section{General setup of LBL experiments and their beam specifications}
In long-baseline experiments, neutrino beams are produced with well-defined energies and a minimal uncertainty in neutrino flux. These beams typically span from a few hundred MeV to a few GeV, while neutrinos traverse hundreds of kilometers from their source to the detector. The process begins with the production and acceleration of high-energy proton beams, reaching tens to hundreds of GeV, which are then directed onto a solid graphite target. As these energetic protons strike the target, they come to a halt, triggering a cascade of pions and kaons, which play a crucial role in neutrino production \cite{Feldman:2012jdx, Diwan:2003bp}.
\par Before decaying, pions ($\pi^\pm$) and kaons ($K^\pm$) pass through magnetic horns (conductors carrying strong pulsed electric currents) \cite{Ichikawa:2012ey, Kahn:2003zz, Baussan:2011yg}. These horns use the Lorentz force, generated by both electric and magnetic fields, to selectively focus either positively or negatively charged particles into a well-defined beam, while deflecting the unwanted ones. For example, when $\pi^+$ and $K^+$ are focused, they produce neutrinos, whereas focusing $\pi^-$ and $K^-$ results in antineutrinos. While the selected particles move forward, others are steered away by the electromagnetic force (Lorentz force) inside the horn. Finally, the focused particles (mostly $\sim99\%$ pions) enter into a decay pipe — an evacuated tunnel about 50 to 200 meters long — where they undergo further decay as per the following decay channels.
\begin{align}
    \pi^+ \rightarrow \mu^++\nu_\mu,\nonumber\\
   \pi^- \rightarrow \mu^-+\bar{\nu}_\mu.
   \label{eq:nu_production_LBL}
\end{align}
The branching ratio of pions decaying into muons is $>99\%$. While pions can also decay through other channels, such as $\pi^+ \rightarrow e^+ + \nu_e$, these occur with a negligible branching ratio of about $10^{-4}$. The low branching ratio for the decay channel of $\pi^+ \rightarrow e^+ + \nu_e$ is due to the helicity suppression. To ensure a pure, nearly mono-energetic neutrino ($\nu_\mu$) beam, a beam dump is used to absorb any remaining particles, including muons, as well as leftover pions and kaons.\\
\subsection{Concept of near detector and far detector in LBL setup}
Long-baseline neutrino oscillation experiments investigate how neutrinos change their flavors as they propagate over distances of hundreds to thousands of kilometers. Because neutrino oscillation probabilities depend on both energy and baseline, these experiments utilize two detectors: a near detector (ND) situated close to the neutrino source and a far detector (FD) placed at a distant location. Comparing measurements from these near and far detectors enables us to determine the fundamental neutrino oscillation parameters while mitigating the systematic uncertainties.\\
\par \textbf{Near detectors} (ND) are kept close to the source of the neutrinos ($\sim$ few hundred meters). Sometimes it consists of several detector components to study different types of neutrino interactions. Near detectors (NDs) \cite{DUNE:2021tad, Emberger:2018pgr, Burgman:2021csn} play a crucial role in characterizing the neutrino beam before oscillations occur. These detectors are designed with specific features to measure the neutrino flux, energy spectra, and interaction cross-sections as precisely as possible. It is smaller in size than far detectors, but large enough to collect significant event statistics. The energy resolution of the near detector is very high ($<100$ MeV), reducing the uncertainty in the measured neutrino spectra. It is also enriched with fine-grained tracking to reconstruct interaction vertices and outgoing particles. It also utilizes shielding and timing cuts to eliminate the beam-induced backgrounds. For T2K, there are two near detectors: ND280 (off-axis, plastic scintillator, magnetized) \cite{Yevarouskaya:2023yth, T2K:2022atj} and INGRID (on-axis, iron scintillator, non-magnetized) \cite{Abe:2011xv}. ND280 measures the neutrino flux and cross-section at the same off-axis angle as Super-Kamiokande, whereas INGRID monitors the neutrino beam direction and intensity. For NO$\nu$A, there is a 300-ton liquid scintillator kept at around 1 km away from the neutrino source as a near detector \cite{ NOvA:2024rov} with good energy resolution and particle-tracking capability. In MINOS, at a distance of 1 km away from the source, a 980-ton magnetized ($\sim$ 0.5 Tesla) iron-scintillator tracking calorimeter was placed as a near detector \cite{MINOS:2008hdf}. In DUNE, there will be a near detector (575 m away from the source) having multiple specifications \cite{DUNE:2021tad}. The DUNE ND system is planned to consist of three major sub-detectors: ND-LAr (Liquid Argon Time Projection Chamber - LArTPC), ND-GAr (High-Pressure Gaseous Argon TPC - HPgTPC, also called ``SAND (System for on-Axis Neutrino Detection)", magnetized), and ND-LAr Off-Axis System (associated with DUNE PRISM (Precision Reaction Independent Spectrum Measurement)) \cite{Hasnip:2023ygr, Hasnip:2025gyi}. It is one of the most advanced near detectors. The ND-LAr Off-Axis System will be able to move off-axis up to $\pm~30$ meters to sample different neutrino fluxes and also can move sideways relative to the beam. It will measure neutrino interactions at different beam angles, providing a way to study the neutrino flux in different energy regions. It will also help to constrain systematic uncertainties in flux modeling. \\
\par \textbf{Far detectors} (FDs) are pivotal in long-baseline neutrino oscillation experiments, positioned hundreds to thousands of kilometers away from the neutrino source, typically an accelerator-based beamline. Their primary aim is to observe the beam neutrinos after oscillation, providing key insights into measuring neutrino mixing parameters, mass ordering, and CP violation. To enhance the interaction rates, far detectors are significantly larger (on the kiloton scale) than near detectors, as neutrinos interact weakly with matter. They are placed deep underground to reduce cosmic background interference, ensuring cleaner signal detection. Additionally, their design enables precise identification of different neutrino flavors, key to identify neutrino oscillation. We elaborate the specifications of far detectors of some past, present, and future long-baseline experiments in the upcoming subsections with more details.\\
\subsection{Beam specifications of LBL experiments}
Neutrino beams are categorized as wide-band beams (WBB) or narrow-band beams (NBB) based on their energy spectra, which are intrinsically linked to the focusing and subsequent decay of secondary pions and kaons generated within the target. Likewise, their classification as on-axis or off-axis beams pertains to the trajectory of the neutrinos relative to the decay pipe's alignment.\\
\par A \textbf{wide-band neutrino beam} exhibits a broad and continuous energy distribution shaped by the momentum spectra of secondary pions and kaons \cite{Barger:2006vy, Papadimitriou:2016ksv}. The peak energy of such a beam is experiment-dependent, but generally falls within a few GeV to approximately 10 GeV. This broad energy spectrum facilitates the detection of multiple oscillation peaks (or a range of \( L/E \) values), enhancing sensitivity to important searches in neutrino oscillation frontiers such as the neutrino mass ordering and CP violation. Additionally, the statistical uncertainty in the beam is comparatively lower. The DUNE experiment plans to utilize a wide-band neutrino beam spanning in the range of 0.1 to 10 GeV, with a peak energy around 2.5 GeV. Similarly, the MINOS experiment utilized a wide-band neutrino beam with a peak energy around 3 GeV.\\
\par A \textbf{narrow-Band neutrino beam} is a beam characterized by a narrow energy range, achieved by selecting pions and kaons within a specific momentum range before they undergo decay into neutrinos. This selection results in a beam where most neutrinos possess similar energies, thereby reducing the uncertainties associated with the energy-dependent effects in neutrino interactions, as it effectively probes a fixed \( L/E \) ratio. Such a beam minimizes the background from deep-inelastic neutrino scatterings, thereby enhancing signal purity in oscillation and cross-section measurements \cite{Steinberger:1988zd, NuTeV:1998aks}. Additionally, it significantly mitigates the energy-smearing effects in cross-section measurements. A narrow-band neutrino beam is generated by precisely controlling the momentum of the parent hadrons (primarily pions and kaons) before they decay into neutrinos. A magnetic horn system is employed to select hadrons within a definite momentum window (e.g., 4–6 GeV/c for a \(\sim\) 2 GeV neutrino beam). Further refinement is achieved through hadron spectrometers and collimators, ensuring that only hadrons with the desired momentum reach the decay pipe. As a result, the produced neutrino spectrum remains narrow and well-defined, peaking at a specific energy. The energy of the neutrino beam and that of the parent pions (for both wide-band and narrow-band beams) are related as follows,
\begin{align}
    E_\nu \approx \frac{0.43\times p_\pi}{1+\gamma^2\theta^2},
    \label{eqn:3.4}
\end{align}
where \( p_\pi \) denotes the momentum of the pion, \( \gamma \) represents its Lorentz boost factor, and \( \theta \) denotes the angle between the trajectory of the emitted neutrino and the decay pipe (or the direction of the parent pion). The energy spread in a narrow-band beam is relatively low (\(\sim 10\text{–}20\%\)), but the statistics is significantly low. The currently running LBL experiments T2K and NO$\nu$A utilize this type of neutrino beam and the upcoming Hyper-Kamiokande experiment will also use this kind of neutrino beam from accelerators.\\
\par Neutrino experiments often use high-intensity neutrino beams produced at accelerator facilities. These beams are categorized into on-axis and off-axis configurations depending on the angle at which the neutrinos travel relative to the center of the beamline. The distinction between these two configurations has significant implications for neutrino physics, particularly in long-baseline neutrino oscillation experiments. \\
\par An \textbf{on-axis neutrino beam} consists of neutrinos propagating along the central axis of the hadron decay pipe, meaning their trajectories are collinear with those of the parent mesons \cite{Angelico:2019gyi}. The resulting neutrino energy spectrum is broad, exhibiting a substantial high-energy tail. Since neutrinos originate from the full phase space of the decaying mesons, the neutrino flux is maximized. It is helpful for experiments studying a wide range of neutrino energies, such as cross-section measurements. DUNE and MINOS experiment use the on-axis neutrino beam.\\
\par An \textbf{off-axis neutrino beam} consists of neutrinos that emerge at a small angle \( \theta \) relative to the central axis of the decay pipe \cite{McDonald:2001mc, Karpova:2025lvb}. This is achieved by positioning the detector slightly off (non-collinear) the beam axis. Compared to an on-axis beam, the neutrino energy spectrum is narrower and peaks at a lower energy. Additionally, the suppression of the high-energy tail reduces backgrounds from deep inelastic scattering. The mean neutrino energy is determined by the energy of the parent mesons and their decay kinematics, following the relation in equation \ref{eqn:3.4}. This beam configuration allows for better control over systematic uncertainties in oscillation experiments and helps mitigate backgrounds from neutral-current interactions. Off-axis beams are particularly advantageous in long-baseline neutrino oscillation experiments, as they concentrate the neutrino flux within a narrow energy range where the oscillation probability is maximized, thereby reducing unwanted backgrounds. Most modern long-baseline oscillation experiments, such as T2K (with an off-axis angle of \(2.5^\circ\)) \cite{T2K:2012qoq} and NO$\nu$A (with an off-axis angle of \(0.8^\circ\)) \cite{Lackey:2019niq}, employ an off-axis configuration to enhance the signal over background ratio. By utilizing an off-axis beam, experiments can fine-tune the neutrino energy spectrum to enhance the flux and maximize the oscillation probability.\\
%%%%%%%%%%%%%%%%%%%%%%%%%%%%%%%%%%%%%%%%%%%%%%%%%%%%%%%%
%%%%%%%%%%%%%%%%%%%%%%%%%%%%%%%%%%%%%%%%%%%%%%%%%%%%%%%
\section{A Panorama of LBL Neutrino Experiments}
\subsection{K2K (KEK to Kamioka) [1999 - 2004]}
The K2K (KEK to Kamioka) experiment \cite{K2K:2004iot} was the first long-baseline neutrino oscillation experiment, conducted between 1999 and 2004, aimed at verifying the atmospheric neutrino oscillation results observed by Super-Kamiokande (Super-K) using a controlled neutrino beam. With a 250 km baseline \cite{K2K:2002icj}, K2K employed the 12 GeV Proton Synchrotron (PS) at KEK in Tsukuba, Japan, as its neutrino source. The experimental setup comprised of a near detector at KEK and a far detector at Super-Kamiokande in Kamioka, Japan. A wide-band neutrino beam peaking at 1.3 GeV was generated by bombarding an aluminum target with $6.0\times10^{12}$ protons per ejection, producing pions and kaons, which subsequently decayed into neutrinos. The far detector, a 50 kt water Cherenkov detector, recorded neutrino interactions, while the near detector, with a 1 kt volume, characterized the initial neutrino flux. K2K \cite{Nishikawa:2004qy} successfully observed the disappearance of muon neutrinos, reinforcing the evidence for neutrino oscillations, and measured the mass-squared difference as $\Delta m^2_{32} = (2.8 \pm 0.3) \times 10^{-3}$ eV², while establishing a lower bound on the atmospheric mixing angle, $\sin^2 2\theta_{23} > 0.9$.\\
\subsection{MINOS (Main Injector Neutrino Oscillation Search) [2005 - 2012]/ MINOS+ [2013 - 2016]}
MINOS was a pioneering long-baseline neutrino oscillation experiment with a baseline of 735 km, linking Fermilab in Illinois to the Soudan Mine in Minnesota \cite{MINOS:2014rjg}. It harnessed the high-intensity NuMI (Neutrinos at the Main Injector) beam, a formidable muon neutrino source generated at Fermilab \cite{MINOS:2020llm}. In this setup, 120 GeV protons from the Main Injector bombarded on a graphite target, initiating the production of pions and kaons, which subsequently decayed mostly into muon neutrinos. The experiment’s 5.4 kt far detector, positioned 710 m underground, was meticulously engineered with alternating layers of 2.54 cm thick steel plates and 1 cm thick plastic scintillator strips, enabling precise tracking of neutrino interactions and flavor transitions. MINOS played a pivotal role in refining the atmospheric oscillation parameters and established a stringent lower bound of $\sin^2(2\theta_{23}) > 0.95$.  Over its operational tenure, it collected $1.07 \times 10^{21}$ protons on target (P.O.T) in $\nu_{\mu}$ mode and $0.33 \times 10^{21}$ P.O.T in $\bar{\nu}_{\mu}$ mode. In 2013, MINOS evolved into its successor, MINOS+, which leveraged an upgraded proton beam, accumulating $0.58 \times 10^{21}$ POT in $\nu_{\mu}$ mode within its first two years of operation \cite{Evans:2017brt}. Beyond its precision measurements of atmospheric oscillation parameters, MINOS also offered compelling evidence supporting the non-zero value of $\theta_{13}$ \cite{MINOS:2011amj}.\\
\subsection{OPERA (Oscillation Project with Emulsion-tRacking Apparatus) [2006 - 2018]}
The OPERA experiment \cite{OPERA:2015wbl} was one of the first generation long-baseline neutrino experiments designed to detect the third-generation neutrino, $\nu_\tau$, within a $\nu_\mu$ beam, making $\nu_\mu \rightarrow \nu_\tau$ its primary oscillation channel. Due to its relatively large rest mass energy ($\sim$ 1.7 GeV), the tau lepton ($\tau$) - and consequently $\nu_\tau$ - is challenging to observe in most of the neutrino experiments. Unlike some long-baseline experiments such as MINOS or NO$\nu$A, OPERA primarily relied on a far detector at LNGS without a dedicated near detector \cite{OPERA:2018nar}. However, the CERN Neutrinos to Gran Sasso (CNGS) beamline included monitoring systems to characterize the neutrino beam at its source, effectively serving as a near-detector substitute by analyzing the extracted 400 GeV proton beam from the CERN Super Proton Synchrotron (SPS) before its journey to Gran Sasso. The far detector, located deep underground at the Gran Sasso National Laboratory (LNGS) in Italy, was a hybrid setup combining electronic trackers with emulsion film technology to precisely identify $\nu_\tau$ interactions. With a 730 km baseline and a far detector with a 1.25 kt volume, OPERA employed scintillator strips for event localization. A key component of the detector was the Emulsion Cloud Chamber bricks, which consisted of stacked lead plates interleaved with nuclear emulsion films, enabling high-precision vertex reconstruction. Additionally, resistive plate chambers (RPCs) provided crucial information on tracking and time-of-flight measurements. OPERA utilized a $\sim$17 GeV neutrino beam to effectively detect tau neutrinos, making it a significant contribution in neutrino physics.\\
\subsection{T2K (Tokai-to-Kamioka) [2010 - till date]}
T2K is one of the pioneering second-generation long-baseline experiments featuring both near and far detectors. It employs two near detectors positioned at 280 meters from the neutrino production target at J-PARC: ND280 \cite{Yevarouskaya:2023yth, T2K:2022atj}, which is $2.5^\circ$ off-axis, and INGRID (Interactive Neutrino GRID) \cite{Abe:2011xv}, which is aligned on-axis. ND280 comprises of multiple sub-systems designed for various measurements, including neutral current background identification, high-precision momentum determination, and particle identification using energy loss gradients (dE/dX) through time projection chambers (TPCs). Additionally, it incorporates a fine-grained detector made of polystyrene scintillator bars serving as an active neutrino interaction target and electromagnetic calorimeters for detecting electrons and photons. On the other hand, INGRID consists of 16 iron-scintillator modules arranged in a cross-shaped configuration, facilitating monitoring of beam stability, direction, and intensity. There are ongoing plans to reduce INGRID's energy threshold from 450 MeV/c to 300 MeV/c \cite{Roth:2024mwe}, which would enhance neutron resolution to $(15 - 30)\%$ \cite{Munteanu:2019llq}. T2K utilizes a beam of muon neutrinos (or antineutrinos) generated at the Japan Proton Accelerator Research Complex (J-PARC) in Tokai and detected at the Super-Kamiokande detector, situated 295 km away in Kamioka. A 515 kW beam of 30 GeV protons is directed onto a graphite target, resulting in the production of charged mesons, which subsequently decay in flight, yielding neutrinos. This process produces a narrow-band neutrino flux with a peak energy around 0.6 GeV. The focusing horns can be tuned to select positively or negatively charged particles, thereby determining whether the far detector receives a neutrino-dominant or antineutrino-dominant beam. Then the beam traverses 295 km baseline to reach Super-Kamiokande. Super-Kamiokande, the far detector of T2K, is a massive water Cherenkov detector located 1,000 meters underground in the Mozumi Mine. It consists of a 50 kt volume of ultra-pure water enclosed within a cylindrical stainless-steel tank measuring 39.3 m in diameter and 41.4 m in height (fiducial volume is 22.5 kt). Inside, 13,000 photomultiplier tubes capture Cherenkov radiation emitted when charged particles exceed the speed of light in water. The resulting Cherenkov ring patterns, which differ for various particle types such as muons and electrons, enable particle identification. The T2K experiment primarily investigates the $\nu_\mu\rightarrow\nu_e$ appearance and $\nu_\mu\rightarrow\nu_\mu$ disappearance channels to explore CP violation in the neutrino sector, resolve the $\theta_{23}$ octant degeneracy, and achieve precise measurements of the 2-3 oscillation parameters. Since its commencement in 2010, T2K has accumulated a total exposure of $3.6\times 10^{21}$ protons-on-target (P.O.T.) by 2020, with (1.97 and $1.63) \times 10^{21}$ P.O.T. recorded in antineutrino and neutrino modes, respectively, at the far detector. Future projections estimate a total exposure of $7.8 \times 10^{21}$ P.O.T. with a 750 kW beam power, evenly distributed between neutrino and antineutrino modes, in accordance with ref. \cite{T2K:2014xyt}. For both appearance and disappearance channels, we consider an uncorrelated $5\%$ systematic uncertainty for signal events and $10\%$ for the background events. For both appearance and disappearance channels, we consider $5\%\,(10\%)$ systematic uncertainty for the signal (background) events. The latest results of T2K can be found in ref. \cite{T2K:2023mcm}.\\
\subsection{NO$\nu$A (NuMI Off-Axis $\nu_e$ Appearance) [2014 - till date]}
NO$\nu$A is a prominent second-generation long-baseline neutrino experiment currently operating in the U.S. It utilizes the neutrino beam from the NuMI (Neutrinos at the Main Injector) \cite{Adamson:2015dkw, NOvA:2004blv, NOvA:2019cyt} facility at Fermilab, which runs at a beam power of 700 kW. The experiment's far detector is positioned at 810 km from the source (in Ash River, Minnesota) at an off-axis angle of $0.8^\circ$, with the neutrino beam having an average energy of approximately 2 GeV. The far detector is a 14 kt fiducial mass tracking calorimeter filled with the liquid scintillator. Additionally, NO$\nu$A features a near detector located at 1 km from the source, sharing the same off-axis angle as the far detector. The primary oscillation channels under investigation in this experiment include $\nu_\mu\rightarrow\nu_e$ and $\bar{\nu}_\mu\rightarrow\bar{\nu}_e$. Like T2K, one of NO$\nu$A’s key objectives is to explore CP violation in the neutrino sector and determine the neutrino mass ordering. Furthermore, it plays a crucial role in improving the precision of neutrino oscillation parameter measurements \cite{NOvA:2021nfi}. In our analysis, we take into account the full projected exposure of NO$\nu$A, totaling $3.6\times10^{21}$ P.O.T, distributed equally between neutrino and antineutrino modes, as outlined in ref. \cite{NOvA:2007rmc}. For both appearance and disappearance channels, we assume systematic uncertainties of $5\%$ for signal events and $10\%$ for the background events.\\
\subsection{T2HK (Tokai-to-Hyper-Kamiokande) experiment in \\JAPAN}
T2HK is an upcoming third-generation long-baseline neutrino experiment designed to advance our understanding of neutrino oscillation parameters. This experiment features a 187 kt water Cherenkov far detector (Hyper-Kamiokande) and will receive beam from the J-PARC facility in Tokai covering a baseline of 295 km. T2HK will employ two near detectors: ND280 (Near Detector 280) and IWCD (Intermediate Water Cherenkov Detector). ND280, positioned approximately 280 m from the beam source, is responsible for measuring the initial neutrino flux and energy spectra before oscillation. It includes improved tracking and calorimetry compared to T2K’s ND280, which helps in reducing systematic uncertainties. IWCD \cite{Akutsu2021IWCD, IWCD_TDR_2022}, a 1 kt detector located at 1 km from the source, enables neutrino flux measurements at varying off-axis angles. The experiment utilizes a 30 GeV proton beam with a beam power of 1.3 MW directed onto a graphite target, generating an artificial, narrow-band neutrino beam with a $2.5^\circ$ off-axis angle. The neutrino beam will achieve its first oscillation maximum at 0.6 GeV and will have an energy range of 100 MeV to 3 GeV.
Unlike DUNE, T2HK operates with an asymmetric runtime distribution: 2.5 years in neutrino mode and 7.5 years in antineutrino mode. This corresponds to the total exposure of $2.7\times10^{22}$ protons-on-target, leading to a total of 2431 kt$\cdot$MW$\cdot$year. The experiment primarily investigates four oscillation processes: $\nu_\mu\rightarrow\nu_e$ appearance, $\bar{\nu}_\mu\rightarrow\bar{\nu}_e$ appearance, $\nu_\mu\rightarrow\nu_\mu$ disappearance, and $\bar{\nu}_\mu\rightarrow\bar{\nu}_\mu$ disappearance. It assumes systematic uncertainties of $5\%$ for appearance channels (with an improved uncertainty of $2.7\%$) and $3.5\%$ for disappearance channels \cite{Hyper-Kamiokande:2018ofw}. T2HK is expected to provide a robust signal of CP violation, resolve the octant degeneracy of $\theta_{23}$.\\

\par The T2HKK (Tokai-to-Hyper-Kamiokande to Korea) experimental setup comprises of two far detectors: one situated in Japan, termed as the Japanese Detector (JD), and the other proposed in Korea, referred to as the Korean Detector (KD). The KD is designed to receive the same neutrino beam as JD, but at an extended baseline of 1100 km, optimizing its sensitivity at the second oscillation maximum. Like JD, the proposed KD will also be a 187 kt water Cherenkov detector, ensuring sufficient statistics at the second oscillation maximum. The expected systematic uncertainties for the JD+KD setup are 5\% for the appearance channel and 3.5\% for the disappearance channel \cite{Panda:2022vdw}.\\
\subsection{DUNE (Deep Underground Neutrino Experiment) in U.S.A}
The upcoming international DUNE project hosted by Fermilab represents a major step forward in the roadmap of neutrino physics. It is another important next-generation long-baseline experiment planned to begin data collection around 2030. One of its main goals is to study CP violation in the neutrino sector, which could help to explain why the Universe has more matter than antimatter. DUNE is expected to resolve the octant ambiguity of $\theta_{23}$ and determine the correct order of neutrino masses due to significant Earth's matter effect that neutrinos feel while traveling from Fermilab to Homestake. The experiment \cite{DUNE:2018tke} has a long 1285 km baseline, stretching from Fermilab to the Sanford Underground Research Facility (SURF) in South Dakota. The far detector is a 40 kt liquid argon time projection chamber (LArTPC). The experiment is planned to be upgraded gradually over time. The neutrino beam is produced by the Main Injector at the Long-Baseline Neutrino Facility (LBNF), where a 1.2 MW, 120 GeV proton beam strikes a graphite target, creating charged mesons that decay into neutrinos. This results in a wide-band neutrino flux, spanning from a few hundred MeV to tens of GeV, peaking at around 2.5 GeV, with most of the neutrinos in the 1 to 5 GeV range. The near detector complex consists of three main parts: ND-LAr (Liquid Argon Near Detector) – ArgonCube, ND-GAr (Gaseous Argon Detector) – High-Pressure Gaseous TPC, and SAND (System for on-Axis Neutrino Detection). These components help in calibrating the neutrino flux, reducing uncertainties in cross-section measurements, and checking the stability of the beam over the time. According to the DUNE Technical Design Report \cite{DUNE:2020jqi}, the experiment will run for a total of 10 years, with equal time spent in neutrino and antineutrino modes (5 years each). This corresponds to an annual exposure of $1.1\times10^{21}$ P.O.T, leading to a total exposure of 480 kt$\cdot$MW$\cdot$year. DUNE will mainly study four key neutrino oscillation channels: $\nu_\mu\rightarrow\nu_e$ appearance, $\bar{\nu}_\mu\rightarrow\bar{\nu}_e$ appearance, $\nu_\mu\rightarrow\nu_\mu$ disappearance, and $\bar{\nu}_\mu\rightarrow\bar{\nu}_\mu$ disappearance, assuming systematic uncertainties of $2\%$ for appearance channels and $5\%$ for disappearance channels.











 
 

















































%%%%%%%%%%%%%%%%%%%%%%%%%%%%%%%%%%%%%%%%%%%%%%%%%%%%%%%%%
%%%%%%%%%%%%%%%%%%%%%%%%%%%%%%%%%%%%%%%%%%%%%%%%%%%%%%%








%%%%%%%%%%%%%%%%%%%%%%%%%%%%%%%%%%%%%%%%%%%%%%%%%%%%%%%%%%%%%